%====================================================================
%  高中数学好题积累
%  使用方法：遇到好题时，按专题追加到本文件中。
%  答案开关：\showanswerstrue 显示答案；\showanswersfalse 隐藏答案。
%====================================================================
\documentclass[zihao=-4]{exam-zh}

\usepackage{amsmath}
\geometry{a4paper,margin=2cm}
\setchoices{max-columns=4}
\questionsetup{show-points=true}
\fillinsetup{no-answer-type=none}

%==================== 答案显示开关 ====================
\newif\ifshowanswers
\showanswerstrue
% \showanswersfalse

\ifshowanswers
  \examsetup{
    question/show-answer=true,
    solution/show-solution=show-stay,
  }
\else
  \examsetup{
    question/show-answer=false,
  }
\fi

%==================== 新题模板 ====================
% \section{专题名称}
% \begin{question}[points=5]
% 题目正文 \paren[A]
% \begin{choices}
%   \item 选项 A
%   \item 选项 B
%   \item 选项 C
%   \item 选项 D
% \end{choices}
% \begin{solution}
% 解析与方法小结。
% \end{solution}
% \end{question}

\begin{document}

\begin{center}
  {\zihao{3}\bfseries 高中数学好题积累}\\[4pt]
  {\zihao{4}\bfseries 按专题持续整理}
\end{center}

\noindent\textbf{说明：}本文件是长期维护的好题总库。题目按知识专题分类，每题尽量保留答案、解析和方法小结。

\section{排列组合与拉丁方}
	
	甲，乙、丙三人未来3 周均要去 A、B、C 三个地方调研，每人每个地方调研时长	为1周，如每个人随机安排顺序，则每周三个人去的地方都不同的概率为 \textbf{$\frac{1}{18}$}
	
	\href{https://en.wikipedia.org/wiki/Latin_square}{拉丁方 Latin Square}


\section{数列·单项选择题}

\begin{question}[points=5]
一个数列的前几项为 $1,3,6,10,15,21,28,36,\ldots$，从第二项起，每一项与前一项的差依次构成等差数列，则该数列的第 $50$ 项为 \paren[A]
\begin{choices}
  \item $1275$
  \item $1596$
  \item $1597$
  \item $1598$
\end{choices}
\begin{solution}
相邻两项之差依次为 $2,3,4,\ldots$，故
\[
  a_n=1+(2+3+\cdots+n)=1+\frac{(n-1)(n+2)}2=\frac{n(n+1)}2.
\]
所以 $a_{50}=\dfrac{50\times51}{2}=1275$，故选 A。
\end{solution}
\end{question}

\begin{question}[points=5]
已知数列 $\{a_n\}$ 满足
\[
  (n-1)a_n=(n+1)a_{n-1}\quad(n\ge 2),\qquad a_1=2,
\]
则满足 $a_n<930$ 的最大正整数 $n$ 为 \paren[B]
\begin{choices}
  \item $28$
  \item $29$
  \item $30$
  \item $31$
\end{choices}
\begin{solution}
由递推式得
\[
 \frac{a_n}{a_{n-1}}=\frac{n+1}{n-1},
\]
于是
\[
 a_n=2\cdot\frac31\cdot\frac42\cdots\frac{n+1}{n-1}=n(n+1).
\]
因为 $29\times30=870<930$，而 $30\times31=930$，所以最大正整数为 $29$，故选 B。
\end{solution}
\end{question}

\begin{question}[points=5]
已知数列 $\{a_n\}$ 满足 $a_1=4$，
\[
 a_{n+1}=4a_n-12n+4,
\]
则 $a_n=$ \paren[D]
\begin{choices}
  \item $n$
  \item $2n$
  \item $3n$
  \item $4n$
\end{choices}
\begin{solution}
设 $b_n=a_n-4n$，则
\[
 b_{n+1}=a_{n+1}-4(n+1)=4a_n-16n=4b_n.
\]
又 $b_1=a_1-4=0$，所以 $b_n=0$，即 $a_n=4n$，故选 D。
\end{solution}
\end{question}

\begin{question}[points=5]
在数列 $\{a_n\}$ 中，$a_1=1$，且 $a_n+a_{n+1}=2n$，则 $\{a_n\}$ 的通项公式为 \paren[B]
\begin{choices}
  \item $a_n=n$
  \item $a_n=\begin{cases}n,&n\text{ 为奇数},\\n-1,&n\text{ 为偶数}
  \end{cases}$
  \item $a_n=\begin{cases}n,&n\text{ 为奇数},\\n+1,&n\text{ 为偶数}
  \end{cases}$
  \item $a_n=\begin{cases}n-1,&n\text{ 为奇数},\\n+1,&n\text{ 为偶数}
  \end{cases}$
\end{choices}
\begin{solution}
由
\[
 a_n+a_{n+1}=2n,\qquad a_{n+1}+a_{n+2}=2n+2
\]
相减得 $a_{n+2}-a_n=2$。奇数项与偶数项分别成等差数列。又 $a_1=1$，$a_2=1$，故
\[
 a_{2k-1}=2k-1,\qquad a_{2k}=2k-1.
\]
故选 B。
\end{solution}
\end{question}

\begin{question}[points=5]
已知数列 $\{a_n\}$ 的前 $n$ 项和为
\[
 S_n=\frac32n^2-\frac12n,
\]
令 $b_n=\dfrac1{a_na_{n+1}}$，$T_n$ 为数列 $\{b_n\}$ 的前 $n$ 项和，则 $T_n=$ \paren[B]
\begin{choices}
  \item $\dfrac13-\dfrac1{3n+1}$
  \item $\dfrac13-\dfrac1{3(3n+1)}$
  \item $\dfrac13-\dfrac1{3(3n-1)}$
  \item $\dfrac13+\dfrac1{3(3n+1)}$
\end{choices}
\begin{solution}
由 $a_n=S_n-S_{n-1}$（$n=1$ 时也成立），得 $a_n=3n-2$。因此
\[
 b_n=\frac1{(3n-2)(3n+1)}
 =\frac13\left(\frac1{3n-2}-\frac1{3n+1}\right).
\]
裂项相消得
\[
 T_n=\frac13\left(1-\frac1{3n+1}\right)
 =\frac13-\frac1{3(3n+1)}.
\]
故选 B。
\end{solution}
\end{question}

\section{数列·多项选择题}

\begin{question}[points=6]
已知数列 $\{a_n\}$ 满足 $a_1=-2$，
\[
 \frac{a_n}{a_{n-1}}=\frac{2n}{n-1}\quad(n\ge2),
\]
其前 $n$ 项和为 $S_n$，则下列结论正确的是 \paren[ABD]
\begin{choices}
  \item $a_2=-8$
  \item $a_n=-n\cdot2^n$
  \item $S_3=-30$
  \item $S_n=(1-n)2^{n+1}-2$
\end{choices}
\begin{solution}
累乘得
\[
 a_n=-2\prod_{k=2}^{n}\frac{2k}{k-1}=-n2^n,
\]
故 A、B 正确，且 $S_3=-2-8-24=-34$，C 错误。由错位相减公式
\[
 \sum_{k=1}^{n}k2^k=(n-1)2^{n+1}+2,
\]
所以 $S_n=(1-n)2^{n+1}-2$，D 正确。
\end{solution}
\end{question}

\begin{question}[points=6]
数列 $\{a_n\}$ 的前 $n$ 项和为 $S_n$，且 $a_1=1$，
\[
 a_{n+1}=\begin{cases}
 a_n-2^n,&n\text{ 为奇数},\\
 a_n+2^{n+1},&n\text{ 为偶数}.
 \end{cases}
\]
则下列结论正确的是 \paren[ABD]
\begin{choices}
  \item $a_4=-1$
  \item $a_n=\begin{cases}2^n-1,&n\text{ 为奇数},\\-1,&n\text{ 为偶数}
  \end{cases}$
  \item $3S_{2n+1}=2^{2n+1}-6n-2$
  \item 数列 $\{(-1)^na_n\}$ 的前 $2n$ 项和为 $-\dfrac{2(4^n-1)}3$
\end{choices}
\begin{solution}
直接计算得 $a_2=-1$，$a_3=7$，$a_4=-1$。并且
\[
 a_{2k}=-1,\qquad a_{2k+1}=2^{2k+1}-1,
\]
故 A、B 正确。进一步，
\[
 S_{2n+1}=\frac{2^{2n+3}-6n-5}{3},
\]
故 C 错误；而
\[
 \sum_{k=1}^{2n}(-1)^ka_k
 =\sum_{k=1}^{n}\bigl(-a_{2k-1}+a_{2k}\bigr)
 =-2\sum_{k=1}^{n}4^{k-1}
 =-\frac{2(4^n-1)}3,
\]
故 D 正确。
\end{solution}
\end{question}

\section{数列·填空题}

\begin{question}[points=5]
已知数列 $\{a_n\}$ 满足 $a_1=2$，且
\[
 a_{n+1}=a_n^2+4a_n+2,
\]
则 $a_3=$ \fillin[254]。
\begin{solution}
由 $a_{n+1}+2=(a_n+2)^2$，令 $b_n=\log_2(a_n+2)$，则 $b_{n+1}=2b_n$，且 $b_1=2$。因此 $b_n=2^n$，
\[
 a_n=2^{2^n}-2,
\]
所以 $a_3=2^8-2=254$。
\end{solution}
\end{question}

\begin{question}[points=5]
已知正项数列 $\{a_n\}$ 满足 $a_1=1$，$a_na_{n+1}=4^n$，则
\[
 a_n=\fillin[{$\begin{cases}2^{n-1},&n\text{ 为奇数},\\2^n,&n\text{ 为偶数}.
\end{cases}$}]
\]
\begin{solution}
令 $x_n=\log_2a_n$，则 $x_1=0$，且 $x_n+x_{n+1}=2n$。由相邻两式相减得 $x_{n+2}-x_n=2$。奇、偶项分别递推，得到
\[
 x_n=\begin{cases}n-1,&n\text{ 为奇数},\\n,&n\text{ 为偶数}.
 \end{cases}
\]
从而得到所填通项。
\end{solution}
\end{question}

\begin{question}[points=5]
已知正项数列 $\{a_n\}$ 满足 $a_1=3$，$a_na_{n+1}=9\cdot2^{2n-1}$，则 $a_n=$ \fillin[{$3\cdot2^{n-1}$}]。
\begin{solution}
由 $a_na_{n+1}=9\cdot2^{2n-1}$ 与 $a_{n+1}a_{n+2}=9\cdot2^{2n+1}$ 相除，得 $a_{n+2}=4a_n$。又 $a_1=3$，$a_2=6$，故奇、偶子数列均为公比 $4$ 的等比数列，统一写为 $a_n=3\cdot2^{n-1}$。
\end{solution}
\end{question}

\section{数列·解答题}

\begin{question}[points=12]
设数列 $\{a_n\}$ 满足
\[
 a_1=2,\qquad a_{n+1}=a_n+4\cdot3^{2n-1}.
\]
\begin{enumerate}
  \item 求数列 $\{a_n\}$ 的通项公式；
  \item 令 $b_n=na_n$，求数列 $\{b_n\}$ 的前 $n$ 项和 $S_n$。
\end{enumerate}
\begin{solution}
(1) 当 $n\ge2$ 时，
\[
 a_n=2+\sum_{k=1}^{n-1}4\cdot3^{2k-1}
 =2+12\sum_{k=1}^{n-1}9^{k-1}
 =\frac{3^{2n-1}+1}{2}.
\]
该式对 $n=1$ 也成立。

(2) 由 $b_k=\dfrac{k}{2}+\dfrac{k3^{2k-1}}2$，得
\[
 S_n=\frac12\sum_{k=1}^{n}k+\frac12\sum_{k=1}^{n}k3^{2k-1}.
\]
利用
\[
 \sum_{k=1}^{n}k9^k=\frac{9-(n+1)9^{n+1}+n9^{n+2}}{64},
\]
可得
\[
 \boxed{S_n=\frac{n(n+1)}4+\frac{3+(8n-1)3^{2n+1}}{128}}.
\]
\end{solution}
\end{question}

\begin{question}[points=12]
设 $S_n$ 为数列 $\{a_n\}$ 的前 $n$ 项和，已知 $a_2=1$，且
\[
 2S_n=na_n.
\]
\begin{enumerate}
  \item 求 $\{a_n\}$ 的通项公式；
  \item 求数列 $\left\{\dfrac{a_{n+1}}{2^n}\right\}$ 的前 $n$ 项和 $T_n$。
\end{enumerate}
\begin{solution}
(1) 令 $n=1$，得 $a_1=0$。当 $n\ge3$ 时，
\[
 2S_n=na_n,\qquad 2S_{n-1}=(n-1)a_{n-1}.
\]
两式相减得
\[
 (n-2)a_n=(n-1)a_{n-1}.
\]
结合 $a_2=1$，累乘得 $a_n=n-1$（$n\ge2$）。因此
\[
 \boxed{a_n=\begin{cases}0,&n=1,\\n-1,&n\ge2.
 \end{cases}}
\]

(2) $a_{k+1}=k$，所以
\[
 T_n=\sum_{k=1}^{n}\frac{k}{2^k}.
\]
由错位相减可得
\[
 \boxed{T_n=2-\frac{n+2}{2^n}}.
\]
\end{solution}
\end{question}

\begin{question}[points=12]
已知数列 $\{a_n\}$ 满足 $a_1=1$，且
\[
 \frac{a_n}{a_{n+1}}=1+2a_n.
\]
\begin{enumerate}
  \item 求 $\{a_n\}$ 的通项公式；
  \item 设 $c_n=4n^2a_na_{n+1}$，求数列 $\{c_n\}$ 的前 $n$ 项和 $T_n$。
\end{enumerate}
\begin{solution}
(1) 原式可化为
\[
 \frac1{a_{n+1}}=\frac1{a_n}+2.
\]
故 $\left\{\dfrac1{a_n}\right\}$ 是首项为 $1$、公差为 $2$ 的等差数列，于是
\[
 \boxed{a_n=\frac1{2n-1}}.
\]

(2)
\[
 c_n=\frac{4n^2}{(2n-1)(2n+1)}
 =1+\frac1{(2n-1)(2n+1)}
 =1+\frac12\left(\frac1{2n-1}-\frac1{2n+1}\right).
\]
因此
\[
 T_n=n+\frac12\left(1-\frac1{2n+1}\right)
 =\boxed{\frac{2n(n+1)}{2n+1}}.
\]
\end{solution}
\end{question}

\begin{question}[points=12]
已知数列 $\{a_n\}$ 满足
\[
 a_0=\frac12,\qquad a_n=a_{n-1}+\frac1{n^2}a_{n-1}^2\quad(n\in\mathbb N^*).
\]
证明：
\[
 \frac{n+1}{n+2}<a_n<n.
\]
\begin{solution}
当 $n=1$ 时，$a_1=\dfrac34$，显然有 $\dfrac23<a_1<1$。以下设 $n\ge2$。

由递推式知 $a_n>a_{n-1}>0$，且
\[
 \frac1{a_{n-1}}-\frac1{a_n}
 =\frac{a_{n-1}}{n^2a_n}<\frac1{n^2}.
\]
于是
\[
 2-\frac1{a_n}
 <\sum_{k=1}^{n}\frac1{k^2}
 <1+\sum_{k=2}^{n}\frac1{k(k-1)}
 =2-\frac1n,
\]
从而 $a_n<n$。

下面证明左侧不等式。由已经证明的 $a_{n-1}<n-1$，有
\[
 a_n=a_{n-1}\left(1+\frac{a_{n-1}}{n^2}\right)
 <a_{n-1}\frac{n^2+n-1}{n^2}.
\]
故
\[
 \frac1{a_{n-1}}-\frac1{a_n}
 =\frac{a_{n-1}}{n^2a_n}
 >\frac1{n^2+n-1}
 >\frac1n-\frac1{n+1}.
\]
从 $2$ 到 $n$ 求和，得
\[
 \frac1{a_1}-\frac1{a_n}>\frac12-\frac1{n+1}.
\]
由于 $a_1=\dfrac34$，所以
\[
 \frac1{a_n}<\frac56+\frac1{n+1}<1+\frac1{n+1}=\frac{n+2}{n+1}.
\]
因此 $a_n>\dfrac{n+1}{n+2}$。综上，
\[
 \boxed{\dfrac{n+1}{n+2}<a_n<n}.
\]
\end{solution}
\end{question}

\section{平面向量·综合拓展}

\begin{question}
$\triangle ABC$ 是边长为 $2$ 的等边三角形，$P$ 是平面内任意一点。求
\[
  \overrightarrow{PA}\cdot
  \left(\overrightarrow{PB}+\overrightarrow{PC}\right)
\]
的最小值。

\begin{solution}
设 $O$ 为等边三角形 $ABC$ 的中心，记
\[
 \vec a=\overrightarrow{OA},\qquad
 \vec b=\overrightarrow{OB},\qquad
 \vec c=\overrightarrow{OC},\qquad
 \vec s=\overrightarrow{OP}.
\]
由等边三角形的对称性，
\[
 \vec a+\vec b+\vec c=\vec 0.
\]
因此
\[
 \overrightarrow{PA}=\vec a-\vec s,
\]
且
\[
 \overrightarrow{PB}+\overrightarrow{PC}
 =\vec b+\vec c-2\vec s
 =-\vec a-2\vec s.
\]
从而
\begin{align*}
 \overrightarrow{PA}\cdot
 \left(\overrightarrow{PB}+\overrightarrow{PC}\right)
 &= (\vec a-\vec s)\cdot(-\vec a-2\vec s)\\
 &=2\lvert\vec s\rvert^2-\vec a\cdot\vec s-
   \lvert\vec a\rvert^2\\
 &=2\left\lvert\vec s-\frac14\vec a\right\rvert^2
   -\frac98\lvert\vec a\rvert^2.
\end{align*}
当 $\vec s=\dfrac14\vec a$，即
$\overrightarrow{OP}=\dfrac14\overrightarrow{OA}$ 时取等号。

因为等边三角形的边长为 $2$，所以其外接圆半径为
\[
 \lvert\vec a\rvert=\lvert\overrightarrow{OA}\rvert
 =\frac{2\sqrt3}{3}.
\]
故所求最小值为
\[
 -\frac98\lvert\vec a\rvert^2
 =-\frac98\left(\frac{2\sqrt3}{3}\right)^2
 =\boxed{-\frac32}.
\]

\textbf{方法小结：}利用等边三角形中心向量之和为零，将三个动点向量统一化为 $\overrightarrow{OP}$，再通过向量配方求最值。
\end{solution}
\end{question}

\clearpage
\section{集合专题}

\begin{question}[points=5]
（2013·上海）设常数 $a\in\mathbb R$，集合
\[
 A=\{x\mid (x-1)(x-a)\ge 0\},\qquad
 B=\{x\mid x\ge a-1\}.
\]
若 $A\cup B=\mathbb R$，则 $a$ 的取值范围是 \paren[B]
\begin{choices}
  \item $(-\infty,2)$
  \item $(-\infty,2]$
  \item $(2,+\infty)$
  \item $[2,+\infty)$
\end{choices}
\begin{solution}
\textbf{解法一：分类讨论集合的区间形式。}

当 $a<1$ 时，
\[
 A=(-\infty,a]\cup[1,+\infty),\qquad B=[a-1,+\infty).
\]
因为 $a-1<a$，所以 $(-\infty,a]\cup[a-1,+\infty)=\mathbb R$，
从而此时恒有 $A\cup B=\mathbb R$。

当 $a\ge 1$ 时，
\[
 A=(-\infty,1]\cup[a,+\infty),\qquad B=[a-1,+\infty).
\]
要使两个集合的并集不留下空隙，必须且只需
\[
 a-1\le 1,
\]
即 $a\le 2$。结合 $a\ge1$，得 $1\le a\le2$。

综上，$a\le2$，即
\[
 \boxed{a\in(-\infty,2]}.
\]

\medskip
\textbf{解法二：利用补集。}

条件 $A\cup B=\mathbb R$ 等价于
\[
 (\mathbb R\setminus A)\cap(\mathbb R\setminus B)=\varnothing.
\]
其中
\[
 \mathbb R\setminus B=(-\infty,a-1),
\]
而 $\mathbb R\setminus A$ 是两个根 $1,a$ 之间的开区间。
当 $a\le1$ 时，该开区间位于 $a-1$ 的右侧，交集为空；当 $a>1$ 时，
要使 $(1,a)$ 与 $(-\infty,a-1)$ 不相交，必须且只需 $a-1\le1$。
故仍得 $a\le2$，选择 B。
\end{solution}
\end{question}

\begin{question}[points=5]
（2011·天津）已知集合
\[
 A=\{x\in\mathbb R\mid |x+3|+|x-4|\le9\},
\]
\[
 B=\left\{x\in\mathbb R\;\middle|\;
 x=4t+\frac1t-6,\ t\in(0,+\infty)\right\}.
\]
则集合 $A\cap B=$ \fillin[{$[-2,5]$}]。
\begin{solution}
\textbf{解法一：分段求 $A$，基本不等式求 $B$。}

对集合 $A$，以 $-3$ 和 $4$ 为分界点讨论：
\[
 |x+3|+|x-4|=
 \begin{cases}
  1-2x,&x<-3,\\
  7,&-3\le x\le4,\\
  2x-1,&x>4.
 \end{cases}
\]
由其不大于 $9$，可得
\[
 A=[-4,5].
\]
又因为 $t>0$，由基本不等式
\[
 4t+\frac1t\ge2\sqrt{4t\cdot\frac1t}=4,
\]
等号在 $t=\dfrac12$ 时成立，所以
\[
 x=4t+\frac1t-6\ge-2,\qquad B=[-2,+\infty).
\]
因此
\[
 \boxed{A\cap B=[-2,5]}.
\]

\medskip
\textbf{解法二：几何距离与判别式。}

$|x+3|+|x-4|$ 表示数轴上的点 $x$ 到 $-3,4$ 两点的距离之和。
在区间 $[-3,4]$ 上距离之和恒为 $7$；向两侧每移动 $1$ 个单位，
距离之和增加 $2$。由 $7+2d\le9$ 得 $d\le1$，故 $A=[-4,5]$。

对 $B$，将参数方程改写为
\[
 4t^2-(x+6)t+1=0.
\]
它有正根 $t$ 的充要条件为 $x+6\ge4$，即 $x\ge-2$，
所以 $B=[-2,+\infty)$。再次得到 $A\cap B=[-2,5]$。
\end{solution}
\end{question}

\begin{question}[points=5]
已知
\[
 U=A\cup B=\{x\in\mathbb N\mid 0\le x<6\},
 \qquad A\cap(\mathrm C_U B)=\{1,3,5\},
\]
则集合 $B$ 为 \paren[C]
\begin{choices}
  \item $\{2,4\}$
  \item $\{2,4,6\}$
  \item $\{0,2,4\}$
  \item 不确定
\end{choices}
\begin{solution}
\textbf{解法一：集合运算。}

由题意
\[
 U=\{0,1,2,3,4,5\}.
\]
又因为 $U=A\cup B$，所以凡是属于 $U$ 但不属于 $B$ 的元素，
一定属于 $A$。因此
\[
 A\cap(\mathrm C_U B)=\mathrm C_U B.
\]
已知左边为 $\{1,3,5\}$，故
\[
 \mathrm C_U B=\{1,3,5\},
\]
从而
\[
 B=U\setminus\{1,3,5\}=\boxed{\{0,2,4\}}.
\]

\medskip
\textbf{解法二：逐个判断元素。}

$1,3,5$ 属于 $A\cap(\mathrm C_U B)$，所以它们都不属于 $B$。
而 $U=A\cup B$ 中剩下的元素为 $0,2,4$。若其中任意一个不属于 $B$，
它又必须因 $U=A\cup B$ 而属于 $A$，于是会进入
$A\cap(\mathrm C_U B)$，与该集合只有 $1,3,5$ 矛盾。
所以 $0,2,4$ 都属于 $B$，即 $B=\{0,2,4\}$，选择 C。
\end{solution}
\end{question}

\begin{question}[points=5]
（2025 春·江宁区期末）若命题
\[
 \forall x\in\left[\frac12,2\right],\qquad
 x^2-\lambda x+2\ge0
\]
为真命题，则实数 $\lambda$ 的最大值为 \paren[C]
\begin{choices}
  \item $\dfrac52$
  \item $3$
  \item $2\sqrt2$
  \item $2\sqrt3$
\end{choices}
\begin{solution}
\textbf{解法一：分离参数与基本不等式。}

因为 $x\in\left[\dfrac12,2\right]$ 时 $x>0$，所以
\[
 x^2-\lambda x+2\ge0
 \iff \lambda\le x+\frac2x.
\]
要使不等式对区间内所有 $x$ 都成立，应有
\[
 \lambda\le
 \min_{x\in[1/2,2]}\left(x+\frac2x\right).
\]
由基本不等式
\[
 x+\frac2x\ge2\sqrt2,
\]
等号在 $x=\sqrt2$ 时成立；且 $\sqrt2\in[1/2,2]$。
故
\[
 \boxed{\lambda_{\max}=2\sqrt2},
\]
选择 C。\par\noindent
\textbf{解法二：导数法。}令 $f(x)=x+\dfrac2x$（$x>0$），则
$f'(x)=1-\dfrac2{x^2}$。因此 $f$ 在
$\left[\dfrac12,\sqrt2\right]$ 上递减，在 $[\sqrt2,2]$ 上递增，
最小值在 $x=\sqrt2$ 处取得。由 $f(\sqrt2)=2\sqrt2$，仍得
$\lambda\le2\sqrt2$，故 $\lambda$ 的最大值为 $2\sqrt2$。
\end{solution}
\end{question}

\end{document}
